% GENERATED FILE. Edit canonical lessons, then run npm run generate:books.
% canonical-source-hash: fnv1a64:dcb5bce4738bdf42
% canonical-lessons: MR-C479-nadi, MR-C479-nas, MR-C479-ghasa, MR-C479-pesi, MR-C479-kalij

\chapter{Pulse and Breath}
\label{ch:mr-pulse-and-breath}
\hlchaptermodality{479}

\begin{chapteropening}
\textbf{By the end of this chapter:} \emph{I can name the pulse, a vein, the throat, a cell and the heart as the seat of feeling.}

Complete the last lesson of chapter 479: I can name the pulse, a vein, the throat, a cell and the heart as the seat of feeling.
\end{chapteropening}

\section[\texorpdfstring{\mr{नाडी} (nāḍī)}{नाडी (nāḍī)}]{\mr{नाडी} (nāḍī) — the pulse}

\label{lesson:MR-C479-nadi}



\begin{quote}
Before the new one: say the Marathi for the gut, then the Marathi for a joint.
\end{quote}

\subsection*{You'll want to know: \mr{नाडी}}
\textbf{\mr{नाडी}} — \emph{nāḍī} — \textquotedblleft{}the pulse\textquotedblright{}.

\begin{quote}
\mr{दवाखान्यात} \mr{त्यांनी} \mr{माझी} \mr{नाडी} \mr{तपासली}.
\end{quote}

\begin{quote}
\emph{davākhānyāt tyānnī mājhī nāḍī tapāslī.}
\end{quote}

\begin{quote}
At the clinic they took my pulse.
\end{quote}

Feminine.

\mr{नाडी} \mr{तपासणे}, to take the pulse; also a cord.

\begin{cousinweb}
\textbf{\mr{नाडी}} (\emph{nāḍī}) is a tatsama, Sanskrit \emph{nāḍī}, \textquotedblleft{}tube, vein\textquotedblright{}; CDIAL lists Marathi \emph{nāḍī}, \textquotedblleft{}pulse\textquotedblright{} (CDIAL 7047).
\end{cousinweb}

\begin{grammarlens}[title={What you've built}]
One more word for describing what you lived and felt.
\end{grammarlens}

\subsection*{Your turn}
\begin{itemize}
\raggedright
  \item \emph{Say it:} \emph{nāḍī}
  \item \emph{Say it:} \emph{nāḍī}, once more
  \item \emph{Say it:} \emph{davākhānyāt tyānnī mājhī nāḍī tapāslī.} — At the clinic they took my pulse.
  \item \emph{Recall:} say the Marathi for the gut, then the Marathi for a joint, then say \emph{nāḍī} again
\end{itemize}

\begin{tcolorbox}[breakable,colback=teal!4,colframe=teal!35!black,title={Before you move on}]
What does \mr{नाडी} mean? (\textquotedblleft{}The pulse\textquotedblright{}.) Say it once more.
\end{tcolorbox}
\section[\texorpdfstring{\mr{नस} (nas)}{नस (nas)}]{\mr{नस} (nas) — a vein; a nerve}

\label{lesson:MR-C479-nas}



\begin{quote}
Before the new one: say the Marathi for a joint, then the Marathi for the pulse.
\end{quote}

\subsection*{You'll want to know: \mr{नस}}
\textbf{\mr{नस}} — \emph{nas} — \textquotedblleft{}a vein; a nerve\textquotedblright{}.

\begin{quote}
\mr{रक्त} \mr{नसांमधून} \mr{वाहतं}.
\end{quote}

\begin{quote}
\emph{rakta nasāmmadhūn vāhtaṁ.}
\end{quote}

\begin{quote}
Blood flows through the veins.
\end{quote}

Feminine. Plural: \mr{नसा}, \emph{nasā}.

Plural \mr{नसा}. \mr{नस} \mr{दबणे}, a nerve gets pinched.

\begin{cousinweb}
\textbf{\mr{नस}} (\emph{nas}) is a tadbhava: CDIAL lists Marathi \emph{nas}, \textquotedblleft{}vein, nerve, tendon\textquotedblright{}, under Sanskrit \emph{snasā} (CDIAL 13784).
\end{cousinweb}

\begin{grammarlens}[title={What you've built}]
One more word for describing what you lived and felt.
\end{grammarlens}

\subsection*{Your turn}
\begin{itemize}
\raggedright
  \item \emph{Say it:} \emph{nas}
  \item \emph{Say it:} \emph{nas}, once more
  \item \emph{Say it:} \emph{rakta nasāmmadhūn vāhtaṁ.} — Blood flows through the veins.
  \item \emph{Recall:} say the Marathi for a joint, then the Marathi for the pulse, then say \emph{nas} again
\end{itemize}

\begin{tcolorbox}[breakable,colback=teal!4,colframe=teal!35!black,title={Before you move on}]
What does \mr{नस} mean? (\textquotedblleft{}A vein; a nerve\textquotedblright{}.) Say it once more.
\end{tcolorbox}
\section[\texorpdfstring{\mr{घसा} (ghasā)}{घसा (ghasā)}]{\mr{घसा} (ghasā) — the gullet; the throat}

\label{lesson:MR-C479-ghasa}



\begin{quote}
Before the new one: say the Marathi for the pulse, then the Marathi for a vein.
\end{quote}

\subsection*{You'll want to know: \mr{घसा}}
\textbf{\mr{घसा}} — \emph{ghasā} — \textquotedblleft{}the gullet; the throat\textquotedblright{}.

\begin{quote}
\mr{थंड} \mr{पाणी} \mr{प्यायल्यामुळे} \mr{माझा} \mr{घसा} \mr{बसला}.
\end{quote}

\begin{quote}
\emph{thaṇḍ pāṇī pyāylyāmuḷe mājhā ghasā baslā.}
\end{quote}

\begin{quote}
Drinking cold water made me hoarse.
\end{quote}

Masculine.

\mr{घसा} \mr{बसणे}, to go hoarse; beside \mr{गळा}, the neck and throat.

\begin{cousinweb}
\textbf{\mr{घसा}} (\emph{ghasā}) is a Marathi noun whose origin is not settled.
\end{cousinweb}

\begin{grammarlens}[title={What you've built}]
One more word for describing what you lived and felt.
\end{grammarlens}

\subsection*{Your turn}
\begin{itemize}
\raggedright
  \item \emph{Say it:} \emph{ghasā}
  \item \emph{Say it:} \emph{ghasā}, once more
  \item \emph{Say it:} \emph{thaṇḍ pāṇī pyāylyāmuḷe mājhā ghasā baslā.} — Drinking cold water made me hoarse.
  \item \emph{Recall:} say the Marathi for the pulse, then the Marathi for a vein, then say \emph{ghasā} again
\end{itemize}

\begin{tcolorbox}[breakable,colback=teal!4,colframe=teal!35!black,title={Before you move on}]
What does \mr{घसा} mean? (\textquotedblleft{}The gullet; the throat\textquotedblright{}.) Say it once more.
\end{tcolorbox}
\section[\texorpdfstring{\mr{पेशी} (peśī)}{पेशी (peśī)}]{\mr{पेशी} (peśī) — a cell (of the body)}

\label{lesson:MR-C479-pesi}



\begin{quote}
Before the new one: say the Marathi for a vein, then the Marathi for the gullet.
\end{quote}

\subsection*{You'll want to know: \mr{पेशी}}
\textbf{\mr{पेशी}} — \emph{peśī} — \textquotedblleft{}a cell (of the body)\textquotedblright{}.

\begin{quote}
\mr{प्रत्येक} \mr{जीव} \mr{पेशींनी} \mr{बनलेला} \mr{असतो}.
\end{quote}

\begin{quote}
\emph{pratyek jīv peśīnnī banlelā asto.}
\end{quote}

\begin{quote}
Every living thing is made of cells.
\end{quote}

Feminine.

Feminine; plural \mr{पेशी}.

\begin{cousinweb}
\textbf{\mr{पेशी}} (\emph{peśī}) is a tatsama, taken from Sanskrit \emph{peśī}, \textquotedblleft{}piece of flesh, muscle\textquotedblright{} (CDIAL 8384); modern Marathi gave it the sense \textquotedblleft{}cell\textquotedblright{}.
\end{cousinweb}

\begin{grammarlens}[title={What you've built}]
One more word for describing what you lived and felt.
\end{grammarlens}

\subsection*{Your turn}
\begin{itemize}
\raggedright
  \item \emph{Say it:} \emph{peśī}
  \item \emph{Say it:} \emph{peśī}, once more
  \item \emph{Say it:} \emph{pratyek jīv peśīnnī banlelā asto.} — Every living thing is made of cells.
  \item \emph{Recall:} say the Marathi for a vein, then the Marathi for the gullet, then say \emph{peśī} again
\end{itemize}

\begin{tcolorbox}[breakable,colback=teal!4,colframe=teal!35!black,title={Before you move on}]
What does \mr{पेशी} mean? (\textquotedblleft{}A cell (of the body)\textquotedblright{}.) Say it once more.
\end{tcolorbox}
\section[\texorpdfstring{\mr{काळीज} (kāḷīj)}{काळीज (kāḷīj)}]{\mr{काळीज} (kāḷīj) — the innermost heart; the liver}

\label{lesson:MR-C479-kalij}



\begin{quote}
Before the new one: say the Marathi for the gullet, then the Marathi for a cell.
\end{quote}

\subsection*{You'll want to know: \mr{काळीज}}
\textbf{\mr{काळीज}} — \emph{kāḷīj} — \textquotedblleft{}the innermost heart; the liver\textquotedblright{}.

\begin{quote}
\mr{ती} \mr{बातमी} \mr{ऐकून} \mr{माझं} \mr{काळीज} \mr{तुटलं}.
\end{quote}

\begin{quote}
\emph{tī bātmī aikūn mājhaṁ kāḷīj tuṭlaṁ.}
\end{quote}

\begin{quote}
My heart broke when I heard that news.
\end{quote}

Neuter.

Neuter: \mr{माझं} \mr{काळीज}. \mr{काळजाचा} \mr{तुकडा}, a darling. \mr{काळजी}, care, is often linked with it.

\begin{cousinweb}
\textbf{\mr{काळीज}} (\emph{kāḷīj}) is a tadbhava: CDIAL lists Marathi \emph{kāḷīj}, \textquotedblleft{}liver, heart\textquotedblright{}, under Sanskrit \emph{kāleyaka} (CDIAL 3103). \mr{काळजी}, care, is often linked with it, but CDIAL does not list it there.
\end{cousinweb}

\begin{grammarlens}[title={What you've built}]
One more word for describing what you lived and felt.
\end{grammarlens}

\subsection*{Your turn}
\begin{itemize}
\raggedright
  \item \emph{Say it:} \emph{kāḷīj}
  \item \emph{Say it:} \emph{kāḷīj}, once more
  \item \emph{Say it:} \emph{tī bātmī aikūn mājhaṁ kāḷīj tuṭlaṁ.} — My heart broke when I heard that news.
  \item \emph{Recall:} say the Marathi for the gullet, then the Marathi for a cell, then say \emph{kāḷīj} again
\end{itemize}

\begin{tcolorbox}[breakable,colback=teal!4,colframe=teal!35!black,title={Before you move on}]
What does \mr{काळीज} mean? (\textquotedblleft{}The innermost heart; the liver\textquotedblright{}.) Say it once more.
\end{tcolorbox}
